\documentclass[11pt]{amsart}
\usepackage[margin=1in]{geometry}
\usepackage{amsmath,amssymb,mathtools}
\usepackage{graphicx}
\usepackage{enumitem}
\usepackage[colorlinks=true,linkcolor=blue,citecolor=blue,urlcolor=blue]{hyperref}

\newcommand{\E}{\mathbf E}
\newcommand{\Pp}{\mathbf P}
\newcommand{\N}{\mathbb N}
\newcommand{\Z}{\mathbb Z}
\newcommand{\T}{\mathbb T}
\newcommand{\1}{\mathbf 1}
\newcommand{\card}[1]{\lvert #1\rvert}

\newtheorem{theorem}{Theorem}
\newtheorem{lemma}{Lemma}
\theoremstyle{definition}
\newtheorem{definition}{Definition}
\theoremstyle{remark}
\newtheorem{remark}{Remark}

\title{A Computable Derandomization of a Thin-Set Construction}
\author{agent\_lane\_03}
\date{27 August 2026}

\begin{document}
\maketitle

\section*{Status and source question}

\textbf{Status: substantial partial result (full effective
derandomization; the ``natural'' closed-form question remains).}

Neuwirth asks on page 3 of \cite{Neuwirth} whether some ``natural'' set of
integers is a $\Lambda(p)$ set for every finite $p$ and fails the Rosenthal
property.  The preceding sentence emphasizes that the random construction
provides no explicit procedure for building such a set.

\begin{center}
\includegraphics[width=.94\textwidth]{figures/open_problem_crop.png}
\end{center}

The theorem below gives a parameter-free deterministic algorithm for such a
set.  Its output is computable, but its exhaustive finite searches are
enormous and it is not a closed formula.  Since ``natural'' has no formal
definition, the result is deliberately not presented as a complete answer to
that adjective.

\section*{Definitions and source ingredients}

For $E\subseteq\Z$ and $p>0$, recall that $E$ is a $\Lambda(p)$ set if,
for some $0<r<p$, the $L^p$ and $L^r$ norms are equivalent on trigonometric
polynomials with spectrum in $E$.  If $E=\{n_j\}_{j\ge1}$ is ordered
increasingly, it is \emph{equidistributed} if
\[
       \frac1m\sum_{j=1}^m z^{n_j}\longrightarrow0
       \qquad(z\in\T\setminus\{1\}).
\]
Neuwirth records that every weakly equidistributed set fails the Rosenthal
property \cite[p.~3]{Neuwirth}.

We use two further ingredients proved in the source.

\begin{enumerate}[label=(N\arabic*)]
\item\label{N1}
For an integer $s\ge2$, call a finite set $F$ \emph{$s$-independent} if
there is no relation
\[
 \sum_{j=1}^m a_jx_j=0,\qquad
 3\le m\le2s,\quad x_j\in F\text{ distinct},\quad
 a_j\in\Z\setminus\{0\},
\]
with $\sum_j a_j=0$ and $\sum_j|a_j|\le2s$.
If $I_k$ form a Littlewood--Paley partition and $F_k\subseteq I_k$ are
$s$-independent, then $\bigcup_kF_k$ is a $\Lambda(2s)$ set
\cite[Proposition 5.4]{Neuwirth}.

\item\label{N2}
If each point of a finite set of size $q$ is selected independently with
probability $\ell/q$, then
\begin{equation}\label{eq:indprob}
 \Pp\{\text{the selected set is $s$-dependent}\}
 \le D_s\frac{\ell^{2s}}q,
\end{equation}
where one may take, for example,
$D_s=(2s)(4s+1)^{2s}$; see \cite[Lemma 5.5]{Neuwirth} and its proof.
\end{enumerate}

We also use the complex Bernstein concentration inequality in
\cite[Lemma 6.1]{Neuwirth}: if independent centered variables satisfy
$|X_j|\le1$ and $\sum_j\E|X_j|^2\le v$, then
\begin{equation}\label{eq:bernstein}
 \Pp\left\{\left|\sum_jX_j\right|\ge a\right\}
 <4\exp\left(-\frac{a^2}{4(v+a)}\right).
\end{equation}

\section*{The result}

\begin{theorem}\label{thm:main}
There is a parameter-free computable set $E\subseteq\N$ such that
\begin{enumerate}[label=\textup{(\roman*)}]
\item $E$ is a $\Lambda(p)$ set for every $p<\infty$;
\item $E$ is equidistributed;
\item consequently, $E$ is not a Rosenthal set.
\end{enumerate}
In particular, there is an explicit deterministic procedure producing a set
with the two opposed properties in Neuwirth's question.
\end{theorem}

\section*{Idea of the proof}

Use the factorial annuli with endpoints $p_k=2^{k!}$.  In block $k$ we ask
for about $(k+2)!$ points.  Two finite requirements are imposed.  First, the
block must have no short additive relations, with the forbidden length
tending to infinity with $k$.  This makes the union $\Lambda(p)$ for every
$p$ by the source's Littlewood--Paley criterion.  Second, every prefix of the
block must have small centered exponential sums on a sufficiently fine grid
of roots of unity.  Bernstein's inequality shows that the second condition
fails with very small probability; Neuwirth's relation count does the same
for the first.  Hence at least one subset of every block passes both finite
tests, and choosing the lexicographically first one is a deterministic,
terminating algorithm.

The prefix condition is essential.  Endpoint discrepancy alone would not
control averages while the enumeration is partway through a large block.
The adaptive threshold used below yields relative error $O(k^{-1/2})$
uniformly at every cutoff in block $k$.

\section*{Construction}

For $k\ge4$, set
\[
 p_k=2^{k!},\qquad B_k=\{p_{k-1}+1,\ldots,p_k\},\qquad
 q_k=\card{B_k},\qquad \ell_k=(k+2)!,\qquad
 \delta_k=\frac{\ell_k}{q_k},\qquad L_k=k!.
\]
Let
\[
 h_k=\max\left(\{1\}\cup
 \left\{2\le s\le k:
 D_s\frac{\ell_k^{2s}}{q_k}\le\frac14\right\}\right).
\]
Because $\log q_k\asymp k!$ while
$\log(D_s\ell_k^{2s})=O_s(k\log k)$ for each fixed $s$, we have
\begin{equation}\label{eq:hk}
                         h_k\longrightarrow\infty.
\end{equation}

Let $\Omega_k$ be the set of all $(4p_k)$th roots of unity.  For
$A\subseteq B_k$, $1\le r\le q_k$, and $z\in\T$, define
\[
 Z_{A,k,r}(z)=
 \sum_{m=1}^r
 \big(\1_A(p_{k-1}+m)-\delta_k\big)z^{p_{k-1}+m},
 \qquad v_{k,r}=r\delta_k,
\]
and put
\[
 T_{k,r}=32\left(\sqrt{(v_{k,r}+1)L_k}+L_k\right).
\]
Call $A$ \emph{admissible} if
\begin{enumerate}[label=(A\arabic*)]
\item $A$ is $h_k$-independent when $h_k\ge2$;
\item $|Z_{A,k,r}(z)|\le T_{k,r}$ for every
$1\le r\le q_k$ and every $z\in\Omega_k$.
\end{enumerate}
Order subsets of $B_k$ lexicographically by their indicator words, let
$E_k$ be the first admissible subset, and define
\[
                              E=\bigcup_{k\ge4}E_k.
\]

\section*{Proof of Theorem \ref{thm:main}}

\subsection*{1. Every block search terminates}

Select the elements of $B_k$ independently with probability $\delta_k$.
For a fixed prefix $r$ and fixed $z\in\Omega_k$, the summands defining
$Z_{A,k,r}(z)$ are independent, centered, bounded by $1$, and have total
variance at most $v_{k,r}$.  If
$b=\sqrt{(v_{k,r}+1)L_k}+L_k$ and $a=32b$, then
\[
 \frac{a^2}{4(v_{k,r}+a)}\ge4L_k.
\]
Indeed, $b^2\ge v_{k,r}L_k$ and $b\ge L_k$, so
$64b^2\ge L_kv_{k,r}+32L_kb$.  Applying
\eqref{eq:bernstein} and taking a union bound over at most $p_k$ prefixes
and $4p_k$ grid points gives
\begin{align}
 \Pp\{\text{(A2) fails}\}
 &\le16p_k^2e^{-4L_k}
 =16\exp\bigl(-(4-2\log2)L_k\bigr)<\frac14.\label{eq:discprob}
\end{align}
For $h_k\ge2$, \eqref{eq:indprob} and the definition of $h_k$ give
\[
 \Pp\{\text{(A1) fails}\}\le\frac14;
\]
for $h_k=1$ there is no condition.  Thus an admissible subset exists for
every $k\ge4$.  The lexicographic search is finite and terminates.

All its tests are decidable.  Condition (A1) searches finitely many integer
relations.  In (A2), the numbers involved are rational combinations of
roots of unity and square roots of rationals, so comparison can be carried
out exactly in finite algebraic extensions of $\mathbb Q$.  Given
$n\in\N$, one locates its unique block and performs this finite search.
Therefore $E$ is a computable set.

\subsection*{2. The \texorpdfstring{$\Lambda(p)$}{Lambda(p)} properties}

The symmetric factorial annuli
\[
 [-p_k,-p_{k-1})\cup(p_{k-1},p_k]
\]
form the gross Littlewood--Paley partition used in \cite[\S5]{Neuwirth}.
Fix $s\ge2$.  By \eqref{eq:hk}, $h_k\ge s$ for all sufficiently large
$k$, so $E_k$ is $s$-independent in all sufficiently large blocks.
Ingredient \ref{N1} shows that the corresponding tail of $E$ is a
$\Lambda(2s)$ set.  Adding the finitely many earlier points preserves that
property.  As $s$ is arbitrary (and lower exponents follow from the usual
interpolation/inclusion property), $E$ is $\Lambda(p)$ for every finite
$p$.

\subsection*{3. Uniform discrepancy through partial blocks}

A trigonometric polynomial of degree at most $p_k$ has, by Bernstein's
inequality, supremum at most five times its maximum on $\Omega_k$: choose a
nearest grid point, use $\|P'\|_\infty\le p_k\|P\|_\infty$, and note that
the angular mesh error is at most $\pi/(4p_k)$.  Hence (A2) implies
\begin{equation}\label{eq:blockdisc}
 \|Z_{E_k,k,r}\|_\infty
 \le160\left(\sqrt{(v_{k,r}+1)k!}+k!\right)
 \qquad(1\le r\le q_k).
\end{equation}

Define a deterministic weight $\delta_n=\delta_k$ for $n\in B_k$ and zero
before $B_4$.  If $N=p_{k-1}+r\in B_k$, put
\[
 \sigma_N=\sum_{n\le N}\delta_n
 =A_k+v_{k,r},\qquad A_k=\sum_{i=4}^{k-1}\ell_i,
\]
and
\[
 D_N(z)=\sum_{n\le N}(\1_E(n)-\delta_n)z^n.
\]
At the end of block $i$, \eqref{eq:blockdisc} is $O(i\,i!)$; consequently
the sum of all completed-block errors before block $k$ is $O(k!)$.  Also
$A_k\ge\ell_{k-1}=(k+1)!$.  From \eqref{eq:blockdisc}, uniformly for
$0\le v=v_{k,r}\le\ell_k$,
\begin{equation}\label{eq:reldisc}
\begin{aligned}
 \frac{\|D_N\|_\infty}{\sigma_N}
 &\le \frac{O(k!)}{A_k+v}
 +160\frac{\sqrt{(v+1)k!}+k!}{A_k+v}\\
 &=O(k^{-1/2})\longrightarrow0.
\end{aligned}
\end{equation}
For the square-root term, use
$v+1\le2(A_k+v)$ and $A_k\ge(k+1)k!$.

Taking $z=1$ in \eqref{eq:reldisc} gives
\begin{equation}\label{eq:count}
 \card{E\cap[1,N]}=\sigma_N+D_N(1)\sim\sigma_N.
\end{equation}
In particular, $E$ is infinite.

\subsection*{4. Equidistribution and failure of Rosenthal}

The block weights $\delta_k=\ell_k/q_k$ are eventually decreasing.  Fix
$z\in\T\setminus\{1\}$.  The geometric partial sums
$S_M(z)=\sum_{n\le M}z^n$ are bounded.  Abel summation, applied after the
irrelevant finite initial segment, therefore shows that
\begin{equation}\label{eq:abel}
 W_N(z):=\sum_{n\le N}\delta_nz^n=O_z(1).
\end{equation}
Since $\sigma_N\to\infty$, equations \eqref{eq:reldisc},
\eqref{eq:count}, and \eqref{eq:abel} give
\[
 \frac{1}{\card{E\cap[1,N]}}
 \sum_{n\in E\cap[1,N]}z^n
 =\frac{W_N(z)+D_N(z)}{\sigma_N+D_N(1)}\longrightarrow0.
\]
This holds at every cutoff $N$, including cutoffs inside a block.  Thus
$E$ is equidistributed.  Neuwirth's cited Lust--Piquard criterion now
implies that $C_E(\T)$ contains a copy of $c_0$, and $E$ is not a Rosenthal
set.  This proves the theorem. \qed

\section*{Verification notes and limitations}

\begin{itemize}
\item The simultaneous block requirements are compatible because their two
failure probabilities sum to less than $1$, not because either event is
assumed independent of the other.
\item The grid test includes $z=1$, so the same discrepancy estimate controls
the cardinality denominator.  It also includes every prefix in a block;
this prevents an endpoint-only gap.
\item The construction is canonical once the displayed parameters and the
lexicographic order are fixed.  It is computable but not computationally
feasible: block sizes are of order $2^{k!}$.
\item This packet does not claim that a brute-force lexicographic object is a
``natural'' closed-form set.  It fully answers the precise effective
procedure version suggested by the sentence immediately preceding the
source question.
\end{itemize}

\section*{Novelty check and human-review recommendation}

The bounded search on 27 August 2026 covered the four run indexes and web/
arXiv queries combining ``explicit'', ``computable'', ``deterministic'',
``equidistributed'', ``$\Lambda(p)$ for all $p$'', and ``Rosenthal''.  It
found Neuwirth's 1999 article (uploaded as arXiv:2603.28227) and the later
random thin-set constructions \cite{LQR,LQR2}, but no matching computable
derandomization.  This was not an exhaustive citation-graph review; novelty
confidence is therefore medium.

Human review should focus on (a) the exact algebraic decidability statement
for the finite grid test, (b) the factorial estimate in
\eqref{eq:reldisc}, and (c) whether the community would regard effective
derandomization as addressing enough of the word ``natural'' to merit a
stronger status.  The present recommendation is to retain
\texttt{partial}, even though Theorem \ref{thm:main} itself is complete.

\begin{thebibliography}{9}
\bibitem{Neuwirth}
S.~Neuwirth,
\emph{Two random constructions inside lacunary sets},
Ann. Inst. Fourier (Grenoble) \textbf{49} (1999), no.~6, 1853--1867;
arXiv:2603.28227,
\url{https://arxiv.org/abs/2603.28227}.

\bibitem{LQR}
D.~Li, H.~Queff\'elec, and L.~Rodr\'iguez-Piazza,
\emph{Some new thin sets of integers in harmonic analysis},
J. Anal. Math. \textbf{86} (2002), 105--138;
arXiv:0912.4214,
\url{https://arxiv.org/abs/0912.4214}.

\bibitem{LQR2}
D.~Li, H.~Queff\'elec, and L.~Rodr\'iguez-Piazza,
\emph{On some random thin sets of integers},
Illinois J. Math. \textbf{54} (2010), 927--944;
arXiv:0904.2507,
\url{https://arxiv.org/abs/0904.2507}.
\end{thebibliography}

\end{document}
